\documentclass[14pt,letterpaper]{extarticle}
\usepackage[utf8]{inputenc}

% Set margins
% https://tex.stackexchange.com/a/327136
% this should replace use of the 'fullpage' package
\usepackage[
    includehead,
    margin=2.5cm,
    bmargin=2cm,
]{geometry}

\usepackage{amssymb}
\usepackage{amsthm}
\usepackage{verbatim} % This works better than comment for some reason
\usepackage{hyperref}
\usepackage{amsmath}
\usepackage{gensymb}
\usepackage{enumitem}
\usepackage{xfrac} % Fractions - https://tex.stackexchange.com/a/230973
\usepackage{xcolor}

% The magic latex drawing thing
\usepackage{tikz}
\usetikzlibrary{snakes,arrows,shapes}
\usepackage[en-US,showdow]{datetime2}

% Colored hyperlinks:
% https://www.overleaf.com/learn/latex/Hyperlinks
\hypersetup{
    colorlinks=true,
    urlcolor=blue
}

% Setup nice headings
% https://www.overleaf.com/learn/latex/Headers_and_footers
\usepackage{fancyhdr}
\usepackage{titling} % Save title for use in fancyhdr

% Shortcuts
\newcommand*{\bP}{\mathbf{P}}
\newcommand*{\bq}{\mathbf{q}}

% vectors
\newcommand*{\bv}{\mathbf{v}}
\newcommand*{\bu}{\mathbf{u}}
% bold vector zero
\newcommand*{\bvz}{\mathbf{0}}

\newcommand*{\Cov}{\operatorname{Cov}}



% Use black tombstone for QED
\renewcommand{\qedsymbol}{$\blacksquare$}

% Important sets
\newcommand*{\C}{\mathbb{C}}
\newcommand*{\R}{\mathbb{R}}
\newcommand*{\Q}{\mathbb{Q}}
\newcommand*{\Z}{\mathbb{Z}}

%
% symbols
%

% name by meaning, not by character
\newcommand{\grad}{\nabla}


% New command for 'partial QED' using white tombstone
\newcommand*{\partialqed}{$\hfill \square$ \\}

\newcommand*{\todo}{{\color{red}\textbf{TODO:}} }

% Increased spaces spacing
\usepackage{setspace} \onehalfspacing

\title{Math 468 - Final Project}
\author{Nicholas Schlabach (nickninja@arizona.edu)}
\DTMsavedate{titledate}{2026-05-05}
\date{\DTMUsedate{titledate}}

%% Copied from https://tex.stackexchange.com/a/2238
%% Modified by Techcable 9/19/21 to use '[...]' instead of '(...)'
\newenvironment{amatrix}[1]{%
  \left[\begin{array}{@{}*{#1}{c}|c@{}}
}{%
  \end{array}\right]
}

\begin{document}
\pagestyle{fancy}
\fancyhead[L]{Nicholas Schlabach}
\fancyhead[C]{\hspace{5mm}\Big(\thetitle{}\Big)}
\fancyhead[R]{{\DTMsetdatestyle{mmddyy}\DTMusedate{titledate}}, Page \thepage}
\fancyfoot{} % no footer
\maketitle
% maketitle automatically sets pagestyle=plain, override this
\thispagestyle{fancy}

\text{ } \\
I aim to prove the following statement: If $A$ is a matrix with positive real entries, there is a unique eigenvalue of largest magnitude, that eigenvalue is real, and the corresponding eigenvector can be chosen to have strictly positive components.
\section*{Lawler's Proof of Perron-Frobenius}
Define $\bu\ge \bv$ if $\forall i, u_i\ge v_i$ and likewise for the other comparison operators. 
\vspace{5mm} \\
Assume that $A$ is a $n \times n$ matrix with positive real coefficients. 
\vspace{2mm} \\
Lemma (a) claims that if $\bv\ge \bvz$ and $\bv \ne \bvz$ then $A\bv>0$. By def. of matrix multiplication: 
        \[ (A\bv)_i=\sum_{j=1}^n a_{ij}\bv_i \]
        Each term of the sum is nonnegative as $a_{ij}>0$ and $\bv\ge 0$. Since $\bv\ne 0$, at least one term $a_{ij}\bv_i$ is nonzero and so the sum $(A\bv)_i$ is positive. This is true for every $i \in \{1,\dots,n\}$, so in fact every entry of $\bv$ is positive and $\bv>0$. This proves problem (a), henceforth referred to as lemma (a).
\vspace{2mm} \\
If $\bv\ge 0$, define $g(\bv)$ as the largest $\lambda$ s.t. $A\bv\ge \lambda \bv$, so $g(\bv)=\max\{\lambda:A\bv \ge \lambda \bv\}$. Clearly $g(\bv) \ge 0$ as $A\bv \ge 0\bv =\bvz$ (still using $\bv\ge 0$). The function $g(\bv)$ is well-defined since $A\bv$ is bounded and the set $\{\lambda:A\bv \ge \lambda \bv\}\cap \mathbb{R}_{\ge0}$ is closed.
\vspace{2mm} \\
Lemma (b): Assume $\bv\ge 0$, I claim $g(\bv)>0$ and $\forall c>0, g(c\bv)=g(v)$. We already know $g(\bv)\ge 0$ from above. Using lemma (a), we know that $A\bv>0$. It follows that there must be some $\lambda>0$ s.t. $A\bv\ge \lambda\bv$ proving the first part of the lemma. Consider the equation $A(c\bv)\ge \lambda (c\bv)$. Dividing both sides by $c$ gives the equation $A\bv\ge \lambda\bv$ (this is valid since we are really just looking at $n$ equations in $\R$). It follows that the equations $A(c\bv)\ge \lambda (c\bv)$ and $A\bv\ge \lambda \bv$ have the same set of solutions $\lambda$ and so $g(\bv)=g(c\bv)$ for $c>0$.
\vspace{2mm} \\
Let $V=\{\bv:\bv \ge 0,\bv\ne \bvz\}, V_{\text{unit}}=\{\bv:\bv \ge 0, \bv \ne \bvz\}, \alpha=\sup g(V)$. \\
Lemma (c) says that $g(c\bv)=g(\bv)$ so we know that the value $\alpha$ is unaffected by the magnitude of the vectors being considered (assuming those vectors are nonnegative). In other words $\alpha=\sup g(V)=\sup(V_{\text{unit}})$. \\
The set $V_{\text{unit}}$ is closed and bounded, so it is compact. The function $g$ is continuous, as small changes in $\bv$ should lead to small changes in the maximum $\lambda$. It follows that $g$ attains a maximum on the set $V_{\text{unit}}$. As argued above, lemma (c) means that $\sup g(V)=\sup g(V_{\text{unit}})$. Since $g$ attains a maximum on $V_{\text{unit}}$ it attains that same maximum on $V$. This means $\alpha$ is not just the supremum of $g(V)$, but a maximum as well. As a result, there exists at least one vector $\bv_\alpha$ such that each2$g(\bv_\alpha)=\alpha$ and $\bv_\alpha \ge 0$. \\
\textit{NOTE}: While the value $\alpha$ is unique, there will always be multiple $\bv_\alpha$ due to the ability to rescale by a constant.
\vspace{2mm} \\
Lemma (c): I now claim that each $\bv_\alpha$ is an eigenvector of $A$ associated with eigenvalue $\alpha$. 
BWOC, assume that either some $\bv_\alpha$ is not an eigenvector of $A$ or its associated eigenvalue is not $\alpha$. In either case, $A\bv_\alpha \ne \alpha\bv_\alpha$. \\
Now consider the quantity $A[A\bv_\alpha-\alpha\bv_\alpha]$.
It follows from the hypothesis that $A\bv_\alpha - \alpha\bv_\alpha\ne \bvz$.
By choice of $\alpha$, $g(\bv_\alpha)=\alpha$, $\bv_\alpha\ge \bvz$ and $A\bv_\alpha \ge \alpha \bv_\alpha$.
It follows that $\mathbf{s}:=A\bv_\alpha -\alpha\bv_\alpha\ge \bvz$.
Since $\mathbf{s}\ne \bvz$ by hypothesis, it satisfies the requirements of lemma (a), so we have $A\mathbf{s}>\bvz$.
Using the linearity of $A$ we get $A\mathbf{s}=A^2\bv_\alpha-\alpha A\bv_a>\bvz$.
It follows that $A(A\bv_a)>\alpha(A\bv_\alpha)$. 
Since every entry of $A(A\bv_a)$ is strictly larger than every entry of $\alpha(A\bv_a)$, there exists some $\epsilon>0$ so that $A(A\bv_\alpha)\ge (\alpha+\epsilon)(A\bv_a)$.
However, by definition of $g$ and $\alpha$, $\alpha$ is the largest $\lambda$ satisfying $A\bu\ge \lambda\bu$ where $\bu \ge \bvz$ is an arbitrary vector.
So finding a larger value $\alpha+\epsilon$ satisfying the same inequality is a contadiction.
\vspace{2mm} \\
Lemma (d): I now claim that the vector $\bv_\alpha\ge \bvz$ satisfying $g(\bv_\alpha)=\alpha$ is unique, provided that it satisfies the normalization condition $\sum_{i=1}^n (\bv_\alpha)_i=1$. \\
% The _N here denotes "nonnegative".
Define the norm $|\bv|_N=\sum_{i=1}^n \bv_i$ for all nonnegative vectors $\bv\ge \bvz$ (equivalent to Manhattan distance since all entries are nonnegative). This means the normalization condition becomes $|\bv_\alpha|_N=1$. \\ 
BWOC, assume that the normalized $\bv_\alpha$ is not unique, so that there is a $\bv_\alpha'\ne \bv_\alpha$ both satisfying $g(\bv_\alpha')=g(\bv_\alpha)=\alpha$ and $|\bv_\alpha|_N=1=|\bv_\alpha'|_N$.

\end{document}